\chapter{Introduction}
\label{ch:introduction}

\section{Background}

The development of mobile communication networks towards
sixth-generation (6G) systems introduces questions about how
future wireless capabilities could affect cybersecurity.
The International Telecommunication Union (ITU) describes
the framework for International Mobile Telecommunications
for 2030 and beyond (IMT-2030). Its envisioned usage scenarios
include artificial intelligence (AI) and communication,
integrated sensing and communication (ISAC), and massive
communication \citep{itu2023imt2030}.

These capabilities build on an existing foundation of mobile
network security. The selected fifth-generation (5G) security
specification defines mechanisms for authentication, access
security, and protection of subscriber identity
\citep{threegpp2022security5g}. This provides a baseline for
examining which security practices could remain applicable
as network capabilities develop and which practices may
require further consideration.

\section{Motivation and Problem Statement}

Research on 6G identifies security and privacy challenges
associated with proposed network architectures and enabling
technologies \citep{porambage2021roadmap}. For example,
research on federated learning examines both its potential
to support network security and the threats affecting the
learning process itself \citep{dealwis2026federated}.
Similarly, the integration of wireless sensing and
communication raises questions about unauthorized sensing,
access to sensing results, and the protection of sensitive
information \citep{etsi2026isacsecurity}.

For cybersecurity practitioners, these developments raise
questions about how to secure, test, and monitor future
wireless systems. Discussing their potential implications
requires attention to the evidence supporting each claim.
An established security mechanism, an envisioned capability,
and a research proposal provide different grounds for
drawing conclusions about operational requirements.

This project therefore examines how published descriptions
of proposed 6G capabilities relate to existing 5G security
practices. It focuses on identifying potential changes to
practitioner tasks while making the assumptions and
uncertainties behind those implications explicit.

\section{Aim and Research Question}

The aim of this project is to review emerging cybersecurity
challenges associated with proposed 6G capabilities and
examine their potential implications for industry through
a comparison with established 5G security mechanisms.

The project addresses the following research question:

\begin{quote}
\emph{How might proposed 6G capabilities change the way
cybersecurity teams secure, test, and monitor wireless
networks compared with 5G?}
\end{quote}

\section{Scope and Intended Contribution}

The review covers three areas: AI in network operations,
integrated sensing and communication, and large-scale
Internet of Things (IoT) connectivity. It draws on
peer-reviewed research and technical documents from
standardization organizations.

The intended contribution is a comparison matrix linking
existing 5G security practices to proposed 6G capabilities,
documented security concerns, and potential implications
for cybersecurity practitioners. Each comparison will
identify its supporting sources and distinguish published
findings from the authors' interpretations. The resulting
analysis will also identify questions that require further
evidence from future implementations.